\documentclass[11pt,a4paper]{article}
\usepackage[T1]{fontenc}
\usepackage[utf8]{inputenc}
\usepackage{lmodern,amsmath,amssymb}
\usepackage[margin=25mm]{geometry}
\usepackage[hidelinks]{hyperref}
\hypersetup{pdftitle={The cost of a mark: Newton's vanishing sagitta and a floor of order $\hbar$},pdfauthor={navstokgap research working notes}}
\newcommand{\tightlist}{\setlength{\itemsep}{0pt}\setlength{\parskip}{0pt}}
\setlength{\emergencystretch}{3em}
\setcounter{secnumdepth}{0}
\begin{document}
\hypertarget{the-cost-of-a-mark-newtons-vanishing-sagitta-and-a-floor-of-order-hbar}{%
\section{\texorpdfstring{The cost of a mark: Newton's vanishing sagitta
and a floor of order
\(\hbar\)}{The cost of a mark: Newton's vanishing sagitta and a floor of order \textbackslash hbar}}\label{the-cost-of-a-mark-newtons-vanishing-sagitta-and-a-floor-of-order-hbar}}

\textbf{Abstract.} Newton reads a force off a trajectory by letting the
sagitta and the area between the inertial tangent and the curve go to
zero, keeping their ratio to the time. Once the comparison must be
\emph{recorded}, it has a floor. For every protocol of marks, made by
any instruments at any times, deciding at error probability \(\epsilon\)
between free motion and a constant force \(F\) over a duration \(\tau\),
whatever the body's initial state under each hypothesis, requires

\[\frac s8\sum_j\Delta(\hat D_j)+\frac J2\sum_j\Delta(\hat X_j)\ \ge\ \hbar\arcsin(1-2\epsilon),
\qquad s=\frac{F\tau^2}{2m},\quad J=F\tau,\]

where \(\hat D_j\) and \(\hat X_j\) are the momentum and position
disturbances of mark \(j\). Newton's sagitta is paired with the impulse
a record leaves undetermined, and his impulse with the displacement. For
marks with uncorrelated Gaussian probes the floor takes the sharp form
\(\tau\Delta E=F^2\tau^3/2m\ge24z_{1-\epsilon}^2\hbar\), about
\(65\hbar\) at five per cent error, because a mark's error and the
impulse it delivers are canonically conjugate; insertion of marks finer
than \(\tau_*=(48z_{1-\epsilon}^2m\hbar/F^2)^{1/3}\) then records free
motion only. Correlated and grid-state probes lower that form, and the
disturbance bound holds for them. The part of the signal that no
preparation of the body can move is a phase: Newton's polygon inscribed
in the parabola agrees with it classically at every vertex and differs
quantum mechanically by \(F^2\sum_j\tau_j^3/(24m\hbar)\), the parabolic
segments of the chords times \(F/2v\hbar\). On the historical side,
Newton's \emph{Opticks} contains a measured least length, the
\(1/89000\)th part of an inch, a posited least impulse, and a
disposition whose period times the corpuscle's momentum is invariant
under refraction. The proposition that would join them to a floor is
that a record leaves an undetermined impulse, and Newton asserted its
negation in print.

Draft, 2026-09-17; revised 2026-09-23. Synthesis of the
\href{newton-mark-floor.md}{mark-floor},
\href{planck-gap-derivation.md}{derivation},
\href{planck-gap-probabilistic.md}{probabilistic},
\href{mark-cost-and-statistical-floor.md}{mark-cost},
\href{record-costs-recoil.md}{recoil},
\href{additive-noise-marks.md}{additive-noise},
\href{record-costs-disturbance.md}{disturbance},
\href{record-distance-path-length.md}{path-length} and
\href{polygon-lift-phase.md}{polygon-lift} notes, which hold the proofs
in full. An adversarial review of the recoil, additive-noise and
disturbance notes on 2026-09-23 found no false result and three proof
gaps, now repaired. The revision of 2026-09-22 restated two floors: the
Gaussian floor needs uncorrelated Gaussian probes, and the aperture
floor takes its area form only at the balanced aperture. Two obligations
remain before submission and are stated in §11. Exploratory; no ledger
promotion.

\hypertarget{the-question}{%
\subsection{1. The question}\label{the-question}}

Galileo's comparison is an inertial horizontal line against a falling
parabola. With mass \(m\), transverse force \(F\), horizontal speed
\(v\) and duration \(\tau\), the two paths are
\(q_{\mathrm I}(t)=(vt,0)\) and \(q_{\mathrm F}(t)=(vt,Ft^2/2m)\), and
the region between them has

\[A=\frac{vF\tau^3}{6m},\qquad s(\tau)=\frac{F\tau^2}{2m},\qquad
\Delta E=\frac{F^2\tau^2}{2m},\qquad \frac{3F}{v}A=\tau\Delta E. \tag{1}\]

The sagitta \(s\), the area \(A\) and the product \(\tau\Delta E\) all
scale as powers of \(\tau\) and vanish together. Newton's Lemma X makes
the first displacement quadratic in the time; Lemma XI, Corollaries 4
and 5, make the curved segment one third of the tangent triangle with
both cubic in the time; Proposition VI reads the force off the limit
\(2ms/\tau^2\); and Proposition I builds an orbit as a polygon struck by
impulses at its vertices, the number of vertices ``augmented in
infinitum''. The construction is designed so that the quantities in (1)
disappear and their ratios survive.

Our question is what stops that refinement when the trajectory has to be
recorded rather than contemplated, and whether the stopping point can be
reached from materials Newton possessed.

The answer has five parts. The geometry alone supplies no floor (§2).
Recording supplies one: in a Gaussian model of marks it is set by a
single number \(\kappa\) with the dimensions of action, sharply, in the
combination \(F^2\tau^3/m\) (§3), and quantum kinematics fixes
\(\kappa\ge\hbar/2\) because a mark's error and its recoil are conjugate
(§4). For probes and instruments of every kind the floor survives as a
bound on disturbance, paired with Newton's sagitta and impulse (§5), and
a phase that no preparation of the body can move carries the same
quantity through Newton's inscribed polygon (§6). Newton's optics
supplies both factors of \(\kappa\), a phase whose rate is proportional
to momentum, and the negation of the premise that joins them (§§7--8).
Newton also attempted a preface founded on the ancients, and the ancient
arguments that bear on his own method are the ones he passed over;
stated as premises about records they generate a hierarchy whose first
rung is Zeno's arrow and whose second is the comparison above (§9).

\hypertarget{the-geometry-has-no-floor}{%
\subsection{2. The geometry has no
floor}\label{the-geometry-has-no-floor}}

\textbf{Proposition 1.} For every partition \(0=t_0<\dots<t_n=\tau\)
with steps \(\tau_j\) and mesh \(|\pi|\), the vertex impulse
\(F\tau_j\), the step sagitta \(F\tau_j^2/2m\) and the step area
\(vF\tau_j^3/6m\) tend to zero with the mesh; the areas sum to at most
\(vF\tau|\pi|^2/6m\); and \(2m\,s(\tau_j)/\tau_j^2=F\) exactly at every
step.

The proof is immediate from (1). Berkeley's objection, that evanescent
increments are ``neither finite Quantities nor Quantities infinitely
small, nor yet nothing'', is answered inside the geometry by the limit
ratio, which is what Newton's closing Scholium to Book I Section I
provides. A floor must therefore come from the act of marking the curve,
and the reading of ``inserting a point'' used here is physical: a point
of the trajectory is inserted by making a mark on the body at a time.

\hypertarget{marks-and-the-floor-they-impose}{%
\subsection{3. Marks, and the floor they
impose}\label{marks-and-the-floor-they-impose}}

\textbf{Definitions.} A \textbf{mark} at time \(t_j\) leaves a datum
about the transverse coordinate. Its \textbf{resolution} \(\delta_j\) is
the spread of the resulting position estimate's error; its
\textbf{recoil} \(\Delta_j\) is the spread of the transverse impulse it
delivers to the body. A protocol is a finite set of marks; it
\textbf{decides} the comparison at error probability \(\epsilon\) when
some test on the full record does so with equal priors. A model
satisfies the \textbf{mark trade-off} with constant \(\kappa\) when
every available mark has \(\delta\Delta\ge\kappa\).

\textbf{The statistical model.} Mark \(j\) returns \(R_j=y(t_j)+\xi_j\)
with \(\xi_j\) centred Gaussian of standard deviation \(\delta_j\), and
delivers an impulse \(\iota_j\), centred Gaussian of standard deviation
\(\Delta_j\), independently, so that

\[y(t)=y_0+v_0t+\theta P(t)+\frac1m\sum_j\iota_j(t-t_j)_+,\qquad
P(t)=\frac{Ft^2}{2m},\]

with \(\theta=0\) under the inertial hypothesis and \(\theta=1\) under
the falling one, and \(y_0,v_0\) unknown. Tests are invariant under the
two unknowns, so the statistic is \(u^{\mathsf T}R\) with
\(\sum_iu_i=0\) and \(\sum_iu_it_i=0\), and the optimal invariant test
has error probability \(\Phi(-d/2)\) with
\(d^2=\sup(u^{\mathsf T}P)^2/(u^{\mathsf T}\Sigma u)\) over such \(u\).

\textbf{Theorem 2.} For every protocol satisfying the mark trade-off
with constant \(\kappa\),

\[d^2\ \le\ \frac{F^2\tau^3}{24\,m\kappa},\]

independently of the number of marks and of their times, and the
constant is sharp. Hence deciding at error probability \(\epsilon\)
requires

\[F^2\tau^3\ \ge\ 96\,z_{1-\epsilon}^2\,m\kappa,\qquad
\tau\Delta E\ \ge\ 48\,z_{1-\epsilon}^2\,\kappa,\qquad
A\ \ge\ \frac{16\,z_{1-\epsilon}^2\,v\kappa}{F},\]

with \(z_{1-\epsilon}=\Phi^{-1}(1-\epsilon)\).

\emph{Proof.} The noise enters \(R\) in two ways: \(\xi_j\) affects
\(R_j\) alone, and \(\iota_j\) affects \(R_i\) for \(i>j\) with
coefficient \((t_i-t_j)/m\). Hence

\[u^{\mathsf T}\Sigma u=\sum_j\delta_j^2u_j^2
+\sum_j\frac{\Delta_j^2}{m^2}S_j^2,\qquad S_j=\sum_{i>j}u_i(t_i-t_j),\]

and the arithmetic--geometric mean inequality applied termwise gives
\(u^{\mathsf T}\Sigma u\ge\frac{2\kappa}{m}\sum_j|u_j||S_j|\).

Introduce \(N(w)=\sum_{i:t_i>w}u_i\) and \(T(w)=\int_w^\tau N\). Since
\(\sum_iu_i=0\), \(N\) vanishes for \(w<0\) and for \(w\ge t_k\); \(T\)
is continuous, piecewise linear, zero on \([\tau,\infty)\), and zero on
\((-\infty,0]\) as well, since \(T(0)=\int_0^\tau N=\sum_iu_it_i=0\) by
the invariance under \(v_0\); it has \(T'=-N\) jumping by \(u_j\) at
\(t_j\) and \(T(t_j)=S_j\). Two identities follow. Fubini and one
integration by parts, using \(P'(0)=0\) and \(T(\tau)=0\), give

\[u^{\mathsf T}P=\int_0^\tau P'N=-\int_0^\tau P'T'=\int_0^\tau P''T
=\frac Fm\int_0^\tau T;\]

and since \(T'\) has bounded variation with jumps \(u_j\) and no other
variation, while \(TT'\) vanishes at both ends,

\[\sum_ju_jT(t_j)=\int_{\mathbb R}T\,dT'=-\int_{\mathbb R}T'^2=-E,
\qquad E=\int_0^\tau T'^2 .\]

Therefore \(\sum_j|u_j||S_j|\ge E\) and
\(u^{\mathsf T}\Sigma u\ge2\kappa E/m\). Since \(T\) vanishes at both
ends of \([0,\tau]\), the Poincaré inequality
\((\int_0^\tau T)^2\le\frac{\tau^3}{12}\int_0^\tau T'^2\) applies, with
equality for \(T\propto w(\tau-w)\) (the Euler--Lagrange equation is
\(T''=\) const), and

\[\frac{(u^{\mathsf T}P)^2}{u^{\mathsf T}\Sigma u}
\le\frac{(F/m)^2\frac{\tau^3}{12}E}{2\kappa E/m}=\frac{F^2\tau^3}{24m\kappa}.\]

Finally \(\Phi(-d/2)\le\epsilon\) requires \(d\ge2z_{1-\epsilon}\). For
sharpness, take marks dense in \([0,\tau]\) with weights making \(N\)
approximate \(2w-\tau\), so that \(T\) approximates \(w(\tau-w)\); the
interior weights are negative where \(T\) is positive, so every
\(u_jS_j\) has one sign and \(\sum_j|u_j||S_j|=E\); balance each mark,
\(\delta_j|u_j|=\Delta_j|S_j|/m\), so the arithmetic--geometric step is
an equality; and make the end marks sharp, which costs nothing since
\(S=T=0\) there. Every inequality then becomes an equality in the limit.
\(\square\)

Three features of the proof carry the paper's claim to universality. The
certificate \(u\) ranges over all invariant tests, so no monitoring
scheme is privileged. The function \(T\) encodes the whole protocol in
one object, and the bound depends on it only through the Dirichlet
energy \(E\), which cancels. And the exponent \(3\) on \(\tau\) comes
from the Poincaré constant of the interval, \(\tau^3/12\), which is
where the combination in (1) originates.

\textbf{Corollary 3 (insertion).} Marks inside a window of duration
\(\tau'\), used by themselves, record the force only if
\(F^2\tau'^3\ge96z_{1-\epsilon}^2m\kappa\). So insertion stops at

\[\tau_*=\left(\frac{96\,z_{1-\epsilon}^2\,m\kappa}{F^2}\right)^{1/3},\]

and finer marks are consistent with free motion at the stated
confidence. Newton's polygon may be refined past \(\tau/\tau_*\)
vertices, and its impulses remain in the geometry, though no record
exhibits them.

\textbf{A worst-case counterpart.} If the marks instead confine position
and impulse to intervals with certainty, the same three steps run with a
separating hyperplane in place of the optimal test direction and give
\(F^2\tau^3>9m\kappa\) for every protocol, with two marks sufficing
above \(36m\kappa\). There the initial velocity is bounded by the
preparation's recoil width instead of being free, so \(T\) need not
vanish at \(0\) and the Cauchy--Schwarz step is the one available. That
version has no quantum instance, since certain position confinement
forces full-support momentum by Paley--Wiener; it is recorded here
because it shows that the argument needs only the quadratic structure.

\hypertarget{the-cost-of-a-mark-is-the-probes-uncertainty-product}{%
\subsection{4. The cost of a mark is the probe's uncertainty
product}\label{the-cost-of-a-mark-is-the-probes-uncertainty-product}}

\textbf{Theorem 4.} Let a mark be realized by the impulsive coupling of
the body's position to a probe,
\(U=\exp(-\tfrac{i}{\hbar}\lambda\hat y\hat P_A)\) with
\([\hat Q_A,\hat P_A]=i\hbar\), the pointer \(\hat Q_A\) being read
afterwards. Then

\[U^\dagger\hat Q_AU=\hat Q_A+\lambda\hat y,\qquad
U^\dagger\hat pU=\hat p-\lambda\hat P_A,\]

so the error operator \(\hat N=\hat Q_A/\lambda\) and the delivered
impulse \(\hat D=-\lambda\hat P_A\) satisfy \([\hat N,\hat D]=-i\hbar\)
and

\[\kappa=\delta\Delta=\Delta\hat N\cdot\Delta\hat D
=\Delta\hat Q_A\cdot\Delta\hat P_A\ \ge\ \frac\hbar2\]

for every probe state, pure or mixed.

\emph{Proof.} With \(A=\tfrac{i}{\hbar}\lambda\hat y\hat P_A\) the
commutators \([A,\hat Q_A]=\lambda\hat y\) and
\([A,\hat p]=-\lambda\hat P_A\) are central, so the
Baker--Campbell--Hausdorff series terminates and gives the two displayed
relations, while \(\hat y\) and \(\hat P_A\) are unchanged. Estimating
the body's position by \(\hat Q_A/\lambda\) leaves the error
\(\hat N=\hat Q_A/\lambda\), and
\([\hat N,\hat D]=-[\hat Q_A,\hat P_A]=-i\hbar\). Robertson's inequality
closes it. \(\square\)

The coupling strength cancels, so the cost is a property of the mark
itself, whatever the strength of the coupling. The record and the recoil
are themselves incompatible, since
\([\hat Q_A+\lambda\hat y,-\lambda\hat P_A]=-i\hbar\lambda\), which is
why no reading of the pointer determines the impulse. A deterministic
probe, whose later position records fix its momentum, has \(\kappa=0\)
and no floor; this is the exact point at which classical record models
differ.

For Gaussian probe states whose position and momentum are uncorrelated,
and the linear dynamics of the comparison, the Wigner function is a
probability density evolving classically, so the statistical model of §3
reproduces the quantum experiment exactly. With Theorem 4,

\[\tau\Delta E\ \ge\ 24\,z_{1-\epsilon}^2\,\hbar,\qquad
\tau_*=\left(\frac{48\,z_{1-\epsilon}^2\,m\hbar}{F^2}\right)^{1/3}, \tag{2}\]

about \(65\hbar\) at \(\epsilon=0.05\) and \(24\hbar\) at
\(\epsilon=0.16\).

\textbf{Correlated marks.} A probe whose position and momentum are
correlated delivers an error and a recoil with correlation
\(\rho=-\operatorname{Cov}(\hat Q_A,\hat P_A)/(\Delta\hat Q_A\Delta\hat P_A)\),
and the Robertson--Schrödinger inequality gives
\(\delta^2\Delta^2(1-\rho^2)\ge\hbar^2/4\). The cross terms enter the
proof of Theorem 2 as \(2c_ju_jS_j/m\), and
\(ax^2+2cxy+by^2\ge2(\sqrt{ab}-|c|)|x||y|\) replaces the
arithmetic--geometric step, so every protocol whose marks have
\(|\rho_j|\le\rho\) needs

\[\tau\Delta E\ \ge\ 24\,z_{1-\epsilon}^2\,\hbar\,\sqrt{\frac{1-\rho}{1+\rho}} , \tag{3}\]

sharp by the dense construction of Theorem 2 with every mark at
correlation \(\rho\) of the helpful sign. Three marks at
\(0,\tau/2,\tau\) with the test \(R_1-2R_2+R_3\) come within a factor
\(4/3\), since the outer marks' recoils are invisible to the test and
the middle mark's noise \(-2\xi_2+(\tau/2m)\iota_2\) is one quadrature
of its probe: they decide once
\(\tau\Delta E\ge32z^2\hbar\sqrt{(1-\rho)/(1+\rho)}\), at every force as
\(\rho\to1\). This is the back-action evasion of Yuen's contractive
states (PRL \textbf{51}, 719, 1983). The proofs are in the
\href{mark-cost-and-statistical-floor.md}{mark-cost note}, Theorem C and
Proposition D.

\hypertarget{every-probe-and-every-instrument-the-record-costs-disturbance}{%
\subsection{5. Every probe and every instrument: the record costs
disturbance}\label{every-probe-and-every-instrument-the-record-costs-disturbance}}

The floors of §§3--4 assume Gaussian probes, and the assumption is
needed. A grid state of the Gottesman--Kitaev--Preskill kind with a real
wavefunction has no position--momentum correlation, yet as the middle
probe of the three-mark protocol it turns the statistic into a comb of
spacing \(\sqrt{2\pi\hbar\tau/m}\), and along the standard family of
approximate grid states the teeth narrow, so the error tends to zero at
any force whose signal is not a multiple of that spacing, at the price
of a recoil that diverges. Grid states read both quadratures of a small
displacement at once, the principle of the displacement sensor of
Duivenvoorden, Terhal and Weigand (PRA \textbf{95}, 012305, 2017). What
holds for every probe state is a bound on the recoil, and for every
instrument a bound on disturbance.

\textbf{Theorem 5.} For every protocol of momentum-transfer marks with
their pointers read and probes in arbitrary states, deciding at error
\(\epsilon\) whatever the body's initial state under each hypothesis
requires
\(\arcsin(1-2\epsilon)\le\hbar^{-1}\min_{a,b}\sum_j\Delta_j|P(t_j)-a-bt_j|\),
with \(\Delta_j\) the recoil spread of mark \(j\). For the constant
force,

\[s\sum_j\Delta_j\ \ge\ 8\hbar\arcsin(1-2\epsilon),\qquad s=\frac{F\tau^2}{2m}, \tag{4}\]

which at certain decision reads \(s\sum_j\Delta_j\ge4\pi\hbar=2h\).

The proof moves the signal into the pointers: translating the body's
initial state along the best straight line \(a+bt\) turns the force into
a translation of each pointer by \(\lambda_jc_j\), with
\(c_j=P(t_j)-a-bt_j\) the line's residual at the mark, and a pointer
registers a translation only through the spread of its momentum, which
is the impulse the body receives. Mandelstam--Tamm and the triangle
inequality over the probes finish it, and running the same chain with
the Bures angle in place of total variation gives the constant
\(\arcsin(1-2\epsilon)\), the sharp single-shot one
(\href{record-distance-path-length.md}{path-length note}, Theorem P). At
fixed \(\epsilon\) the correlated three-mark protocol reaches (4) within
a factor \(\sqrt2z_{1-\epsilon}/\arcsin(1-2\epsilon)\), about \(2.1\) at
five per cent error, at every squeezing. The proofs, the grid
construction and the tightness computation are in the
\href{record-costs-recoil.md}{recoil note}. It is the recoil clause of
M3, and Theorem 6 completes it.

The recoil bound reaches well beyond von Neumann's coupling. For any
coupling and any pointer, the pointer commutes with the body's momentum
after the mark, which gives
\([\hat N,\hat D]=-i\hbar-[\hat y,\hat D]-[\hat N,\hat p]\). So every
mark whose error and impulse are independent of the body has conjugate
error and impulse, which is Ozawa's class of measurements with
independent intervention (PRA \textbf{67}, 042105, 2003; Ann. Phys.
\textbf{311}, 350, 2004); if it also leaves the position alone it is a
von Neumann mark on some canonical pair of the apparatus, a dilation
form of Davies's translation-covariant instruments (1976); and Theorem 5
holds for all such marks, position-displacing ones included
(\href{additive-noise-marks.md}{additive-noise note}). When the noise
depends on the body, the offset between the hypotheses still has to be
carried through each mark, and the cost of carrying it is fixed by what
the mark does to the body.

\textbf{Theorem 6 (every instrument).} Let mark \(j\) be any unitary
coupling of the body to an apparatus followed by reading a pointer, with
momentum and position disturbances
\(\hat D_j=U_j^\dagger\hat pU_j-\hat p\) and
\(\hat X_j=U_j^\dagger\hat yU_j-\hat y\), operators on body and
apparatus that may depend on both. Deciding at error \(\epsilon\)
whatever the body's initial state under each hypothesis requires, for
every line \(a+bt\),
\(\arcsin(1-2\epsilon)\le\hbar^{-1}\sum_j\sup\Delta(c_j\hat D_j-mc_j'\hat X_j)\)
with \(c_j=P(t_j)-a-bt_j\), the supremum over the states that enter the
mark in the interpolating processes of the proof; for the constant
force,

\[\frac s8\sum_j\Delta(\hat D_j)+\frac J2\sum_j\Delta(\hat X_j)\ \ge\ \hbar\arcsin(1-2\epsilon),\qquad J=F\tau. \tag{5}\]

The constants \(1/8\) and \(1/2\) are the price of an initial state that
may differ between the hypotheses: a test allowed to assume the same
initial state under both faces the coefficients \(s\) and \(J\).

The proof is a hybrid argument that carries the offset \((c,mc')\)
between the hypotheses through the marks one at a time. Passing mark
\(j\) applies the unitary \(T^\dagger U_jTU_j^\dagger\), with \(T\) the
translation by the offset; its generator, conjugated back to the state
entering the mark, is \((c_j\hat D_j-mc_j'\hat X_j)/\hbar\), and
Mandelstam--Tamm bounds its effect by that operator's spread. The line
of best uniform approximation gives \(|c_j|\le s/8\) and
\(m|c_j'|\le J/2\) (\href{record-costs-disturbance.md}{disturbance
note}). Von Neumann marks give Theorem 5, and marks that read momentum
pay in position jumps.

Equation (5) involves the disturbances alone, so the error--disturbance
relations, Ozawa's included (PRA \textbf{67}, 042105, 2003), are unused.
Newton's two quantities appear with the partners they have in the
aperture bound of §10: the sagitta with momentum, the impulse with
position.

\hypertarget{the-part-no-preparation-can-move}{%
\subsection{6. The part no preparation can
move}\label{the-part-no-preparation-can-move}}

The floors of §§3--5 are paid in spreads, of resolution, recoil or
disturbance, and a protocol can move a spread from one quantity to
another, as squeezing and grid states do. The comparison also carries a
quantity with no spread at all. Newton's polygon of Proposition I
inscribed in the parabola, with each step's impulse \(F\tau_j\) split
equally between its ends, moves uniformly along each chord and has the
parabola's position and velocity at every vertex, so classically the two
are the same state there.

\textbf{Theorem 7.} For every partition of \([0,\tau]\) into steps
\(\tau_j\), the two evolutions satisfy \(W_{\rm poly}=e^{i\theta_N}W_F\)
with

\[\theta_N=\frac{F^2}{24m\hbar}\sum_j\tau_j^3=\frac{F}{2v\hbar}\sum_jS_j,\]

\(S_j=vF\tau_j^3/12m\) being the parabolic segment cut off by the
\(j\)-th chord. Inserting a vertex at fraction \(\lambda\) of a step
lowers \(\theta_N\) by \(F^2\tau_j^3\lambda(1-\lambda)/(8m\hbar)\), the
inscribed triangle over \(\hbar\), so midpoint halving removes it in the
proportions of Archimedes' exhaustion. The same increments composed in
reverse order differ by the phase
\(\Phi_N=F^2(\tau^3-\sum_j\tau_j^3)/(6m\hbar)\), and
\(\Phi_N+4\theta_N=\tau\Delta E/3\hbar\) for every partition.

The proof uses the Weyl relations alone: an impulse \(J\) at time \(t\)
is a phase-space displacement, two histories with equal endpoints differ
by the enclosed phase-plane area over \(\hbar\) times the identity, and
the per-step lobe is \(F^2\tau_j^3/24m\) (the
\href{polygon-lift-phase.md}{polygon-lift note} has it in full). Since
\(W_{\rm poly}W_F^\dagger\) is a multiple of the identity, no
preparation of the body, squeezed or mixed, changes \(\theta_N\). A
control qubit that selects the history reads it at single-shot error
\(\frac12(1-|\sin(\theta_N/2)|)\), so in one controlled pass a step is
told from its chord at error \(\epsilon\) only if
\(\tau_j\ge(48m\hbar\arcsin(1-2\epsilon)/F^2)^{1/3}\), for every state
of the body; \(k\) passes accumulate \(k\theta_N\) and lower that
threshold by \(k^{1/3}\), so the readout's resource is the number of
passes. The single chord lens and the closed loop of the earlier notes
are instances. In numerical analysis the inscribed polygon is the
kick--drift--kick (Strang) splitting of the propagator, exact for a
linear potential up to this c-number; what is new is its reading in
Lemma XI's areas and Archimedes' segments.

\textbf{Theorem 8 (every force law).} For a force history \(f\) on an
interval of duration \(\tau\), let
\(\mathcal K_\tau[f]=\int\frac12m(\dot y_f-\dot y_{\rm chord})^2dt =\frac1{2m}\iint G_\tau ff\)
be the kinetic action of the motion relative to its chord,
\(G_\tau(s,u)=\min(s,u)(\tau-\max(s,u))/\tau\). Newton's polygon, with
each step's impulse split between its ends so as to match the step's
displacement, differs from the motion by the phase
\(\sum_j\mathcal K_{\tau_j}[f]/\hbar\), which refinement never raises;
and every protocol of uncorrelated Gaussian marks has
\(d^2\le\mathcal K_\tau[f]/\kappa\), sharp for \(f\) of one sign. At
\(\kappa=\hbar/2\) the best recorded deflection is twice the phase
between the motion and its chord. The proofs are in §6 of the
\href{polygon-lift-phase.md}{polygon-lift note}. One functional of the
force history, the part of the motion that Newton's refinement removes,
is converted by \(\hbar\) into a phase and by the mark cost into a
statistical distance; for a constant force it is the chord's segment
times \(F/2v\).

\hypertarget{both-factors-are-in-the-opticks}{%
\subsection{\texorpdfstring{7. Both factors are in the
\emph{Opticks}}{7. Both factors are in the Opticks}}\label{both-factors-are-in-the-opticks}}

Newton's optics supplies the two spreads whose product is \(\kappa\), in
the form of two least quantities of the probe.

\textbf{A least length, measured.} Book II Part III defines the interval
of the fits as ``the space it passes between every return and the next
return'' of the ray's disposition to be reflected, and Proposition XVIII
gives it: for the confine of yellow and orange passing perpendicularly
into air, ``the Intervals of their Fits of easy Reflexion are the
\(1/89000\)th part of an Inch'', that is about \(0.285\,\mu\mathrm{m}\).
A datum obtained through the fits repeats with that period, so it
locates the corpuscle no better; Newton's rings are a gauge of exactly
this kind. This is \(\Delta\hat Q_A\).

\textbf{A least impulse, posited.} Query 29 asks whether the rays of
light are ``very small Bodies emitted from shining Substances'', which
gives each a transverse momentum scale. Newton has no way to measure it.
This is \(\Delta\hat P_A\).

\textbf{Colour dependence, measured.} Observation 13 of Book II Part I
reports the ratio of the red to the violet interval as ``greater than as
3 to 2, and less than as 13 to 8'', and ``By the most of my Observations
it was as 14 to 9''. Query 29 makes the violet corpuscles the least and
the red the biggest. The product \(\Delta\hat Q_A\Delta\hat P_A\) is
therefore colour dependent on Newton-age evidence, and its universality
across the spectrum is what Planck's constant later supplies.

\textbf{Inflexion.} Query 1 asks whether bodies ``act upon Light at a
distance, and by their action bend its Rays'', most strongly at the
least distance, which is the effect that limits a mark's resolution when
the beam is narrowed to improve it.

\textbf{The interval times the momentum is a refraction invariant.} Book
II Part III Prop. XVII makes the intervals in two mediums ``as the Sine
of Incidence to the Sine of Refraction'', from Observation 10 of Part I,
and Prop. X takes light to be ``swifter in Bodies than in Vacuo, in the
proportion of the Sines''. Together they make \(\Lambda v\), and so
\(\Lambda p\), the same in every medium for each colour; that Newton's
interval varies inversely with the corpuscle's speed, as de Broglie's
wavelength does, is an old observation (Sakkopoulos, \emph{Eur. J.
Phys.} \textbf{9}, 123, 1988). A fit counted per interval along the path
therefore counts the Maupertuis action \(\int p\cdot dq\) in units of
\(\Lambda p\), and the inscribed polygon and the parabola of Theorem 7
differ by \((F/v)\sum_jA_j/(2\Lambda p)\) fits, with \(A_j\) Lemma XI's
tangent areas. Since the fit at arrival decides reflection or
transmission (Prop. XII), the two are optically different corpuscles on
Newton's own terms once that count reaches half an interval, told apart
on a single corpuscle, where the quantum phase needs a superposition.
Counting \(2\pi\) per interval, a full return of the disposition, the
fits phase is \(4\theta_N\) at \(\Lambda=\lambda/2\); the wave's own
phase, which advances \(\pi\) per interval, differs by \(2\theta_N\).
Extending Props. X and XVII to a continuously varying speed is ours, and
Prop. XV's oblique rule is more complicated than a count along the path.

All eight passages are held verbatim with line anchors in the
\href{../docs/classics/Newton_Opticks_1730_fits_and_queries.md}{source
companion}, from the fourth edition of 1730.

\hypertarget{the-junction-and-the-proposition-newton-denied}{%
\subsection{8. The junction, and the proposition Newton
denied}\label{the-junction-and-the-proposition-newton-denied}}

Theorems 4 to 6 say what the missing premise is. In the mark-floor note
it was stated as M3, that a mark fixing the position within \(\delta\)
leaves the delivered impulse undetermined within \(\kappa/\delta\).
Theorem 4 identifies \(\kappa\) as the probe's own uncertainty product,
so

\[\text{M3 for Newton's corpuscle}\quad\Longleftrightarrow\quad
\Delta\hat Q_A\cdot\Delta\hat P_A\ \ge\ \text{a positive constant},\]

which is Robertson's inequality stated in Newton's two quantities. The
Gaussian floor also uses a clause M3 carries implicitly, that the
undetermined impulse is unrelated to the mark's error; eq. (3) prices
that clause. Theorems 5 and 6 give M3 its universal form, in which the
resolution drops out: a record of the sagitta leaves the body's impulse
undetermined, a record of the impulse leaves its displacement
undetermined, and the exchange rate is \(\hbar\). With M3, §§3--4 give
the floor and the insertion mesh from Newton's own materials, with
\(\kappa=\Lambda p\) in place of \(\hbar/2\); the identification
\(\Lambda=\lambda/2\), \(p=h/\lambda\) makes \(\Lambda p=h/2\), which
differs from \(\hbar/2\) by \(\pi\).

Newton denied M3 explicitly, in print. Proposition XII of Book II Part
III states that every ray is ``put into a certain transient Constitution
or State, which in the progress of the Ray returns at equal Intervals,
and disposes the Ray at every return to be easily transmitted''. The
disposition is a determinate periodic property carried by the ray, so
its fate at a surface is fixed by its phase, and what a prior record
fails to supply is knowledge rather than determination. On that reading
a probe's position and momentum are both definite, later records can
recover the impulse it delivered, \(\kappa=0\), and the comparison has
no floor. The countermodel is explicit: let the corpuscle travel freely
after the mark to a screen at distance \(D\), and two position records
fix its transverse momentum to within \(p(\delta+\delta_s)/D\), which
vanishes as \(D\) grows.

The historical claim is therefore narrow and checkable. Newton's optics
contains both factors of the product that bounds the recorded
comparison; his mathematics contains the comparison and takes it to
zero; and the one proposition that would join them is an indeterminacy
about the fits whose negation he asserted in print. The gap between the
\emph{Principia}'s vanishing sagitta and the \emph{Opticks}' finite
interval of fits follows from that single commitment, and the separation
of the two books is incidental to it. No counterfactual about what
Newton might have discovered is needed, and none is offered.

Theorem 7 adds a second junction, one with no spread in it. The premise
that Theorem 7 uses is the central extension, that an impulse \(J\) and
a displacement \(d\) compose only up to the phase \(Jd/\hbar\) of their
parallelogram, where Newton's Corollary I composes them exactly.
Newton's fits already carry a phase whose rate is proportional to
momentum, with the refraction invariant \(\Lambda p\) as its unit, and
his determinism in Prop. XII does not touch it. A count of such a phase
is order-sensitive in the parallelogram sense: a displacement \(d\) then
a kick \(J\) counts \(p\,d\), the kick first counts \((p+J)\,d\), and
the difference is the parallelogram's area \(Jd\). So a per-colour form
of the quantum premise that squeezing cannot evade is present in the
\emph{Opticks}; what it lacks is the universality of \(\Lambda p\)
across colours.

\hypertarget{newton-and-the-classics}{%
\subsection{9. Newton and the classics}\label{newton-and-the-classics}}

\hypertarget{what-newton-attempted}{%
\subsubsection{What Newton attempted}\label{what-newton-attempted}}

In the early 1690s Newton drafted a set of Classical Scholia to
Propositions IV to IX of Book III, arguing that the oldest philosophers
had known universal gravitation and its inverse-square law, so that the
\emph{Principia} recovered a lost wisdom rather than announcing a
novelty. He never printed them. He did pass material to David Gregory,
whose \emph{Astronomiae physicae et geometricae elementa} of 1702 opens
with the programmatic sentence that lest the physics delivered here seem
something new and unheard of in astronomy, \emph{eandem vetustissimis
Philosophis notam \ldots{} ostendam}, I shall show it was known to the
most ancient philosophers, naming Anaxagoras, Archelaus, Euripides and
Pythagoras. That preface is the printed witness closest to the
unpublished scholia, and the
\href{../docs/classics/Gregory_AstronomiaeElementa_Praefatio_1702_OCR.md}{source
companion} holds its programmatic paragraph. The scholarship is Casini,
``Newton: The Classical Scholia'', \emph{History of Science} \textbf{22}
(1984), 1--46, and Mosley, ``The Origins and Sources of Newton's
Classical Scholia'', \emph{Erudition and the Republic of Letters}
\textbf{9} (2024), 171--217, both verified at metadata level here and
both still to be read.

So Newton did attempt a preface founded on the classics, and the attempt
was doxographic. He sought ancient authority for a law he had already
proved, and the authority he sought was for the \emph{conclusions} of
Book III.

Book I has one scholium on method, closing Section I, and it does touch
the ancient question, from the geometers' side. Newton prefers limits to
the method of indivisibles, answers the objection that ultimate ratios
would imply ultimate magnitudes, so that quantity would consist of
indivisibles, \emph{contra quam Euclides de incommensurabilibus, in
libro decimo Elementorum, demonstravit}, and concludes \emph{cave
intelligas quantitates magnitudine determinatas, sed cogita semper
diminuendas sine limite}: do not take them as quantities of determinate
magnitude, but think of them as always diminishing without limit
(\href{../docs/Newton_Principia_BookI_SectionI_Motte1729_Wilkins2002.md}{Section
I text}). Euclid is cited against least magnitudes, and the physical
dispute about the cut, with the atomists who held least parts, is left
aside.

\hypertarget{the-ancient-arguments-he-passed-over}{%
\subsubsection{The ancient arguments he passed
over}\label{the-ancient-arguments-he-passed-over}}

The ancient material that bears on the \emph{method} of Book I is the
debate over division and the cut, and it is absent from the scholia. Six
positions, each held here in a primary witness:

\begin{itemize}
\tightlist
\item
  \textbf{Democritus's cone} (Plutarch, \emph{De communibus notitiis}
  39). Cut a cone by a plane parallel to the base: are the surfaces of
  the two adjacent sections equal or unequal? If equal, the cone is a
  cylinder; if unequal, it is stepped. Chrysippus answers that the
  surfaces are neither equal nor unequal while the bodies are unequal,
  and Plutarch calls that a licence to write whatever comes to mind.
\item
  \textbf{Zeno's arrow} (Aristotle, \emph{Physics} VI.9). At every now
  the flying arrow occupies a space equal to itself and so is at rest,
  and what is at rest at every now does not move. Aristotle's diagnosis
  is that the argument assumes time to be composed of nows.
\item
  \textbf{The Mohist endpoint} (\emph{Mozi}, Canons and Explanations).
  Halving terminates at an endpoint that cannot itself be halved.
\item
  \textbf{Liu Hui's cutting} (commentary on the \emph{Nine Chapters},
  263 CE). The finer the cutting the smaller the loss; cut and cut again
  until it cannot be cut, and the polygon coincides with the circle with
  nothing lost.
\item
  \textbf{Hui Shi's stick} (\emph{Zhuangzi} 33). A stick one foot long,
  halved every day, is not exhausted in ten thousand generations.
\item
  \textbf{The leap and the time-atom} (al-Shahrastani on al-Nazzam;
  Maimonides, \emph{Guide} I.73). The interval is divisible, yet it is
  crossed by leaps; and for the kalam, time is composed of atoms.
\item
  \textbf{The Vaiśeṣika arrow} (\emph{Vaiśeṣikasūtra} 5.1.16--18,
  layered, c.~100 BCE--200 CE, attribution to Kaṇāda traditional). The
  arrow's particular conjunctions are \emph{ayugapat}, non-simultaneous,
  and that is the ground of the plurality of its motions; impulsion
  causes only the \textbf{first} motion, and each later one is caused by
  the \emph{saṃskāra} deposited by the motion before it, until the
  impression is spent and gravity takes over.
\item
  \textbf{Sound produced from sound} (\emph{Vaiśeṣikasūtra} 2.2.36--37
  in the Candrānanda recension, 2.2.31 in the vulgate). Sound arises
  from conjunction, from disjunction, and from sound, and is
  non-eternal, so what reaches the ear is a later member of a generated
  chain. Vātsyāyana's \emph{Nyāyabhāṣya}, within the window at about 400
  to 450 CE, argues the point from observation: the axe-blow is still
  heard at a distance after the axe-wood contact has ceased.
\item
  \textbf{No motion at all} (Vasubandhu, \emph{Abhidharmakośabhāṣya} ad
  IV.2, c. 350--450 CE). \emph{na gatir yasmāt saṃskṛtaṃ kṣaṇikam},
  there is no motion, because the conditioned is momentary; everything
  conditioned ``is destroyed in the very place where it arose'', so its
  passage to another place is impossible. Zeno concludes this and treats
  it as absurd; Vasubandhu concludes it and accepts it.
\item
  \textbf{A medium that permits motion} (Umāsvāti,
  \emph{Tattvārthasūtra} 5.17--18,

  \begin{enumerate}
  \def\labelenumi{\alph{enumi}.}
  \setcounter{enumi}{2}
  \tightlist
  \item
    2nd--5th c.~CE). \emph{gatisthityupagrahau dharmādharmayor
    upakāraḥ}: the assisting of motion and of rest is the function of
    \emph{dharma} and \emph{adharma}. These substances push nothing;
    they are the standing condition without which motion could not
    occur, and space is given the separate office of accommodation.
  \end{enumerate}
\item
  \textbf{Sound as spreading waves} (Diogenes Laertius VII.158,
  reporting Stoic doctrine). Hearing occurs when the air between is
  struck, ``a vibration which spreads spherically and then forms waves
  and strikes upon the ears, just as the water in a reservoir forms wavy
  circles when a stone is thrown into it''. The same water-ring image
  carries the Vaiśeṣika \emph{vīcī-santāna}, so the model is attested
  independently in both traditions.
\item
  \textbf{Continuity that fails below sense} (Epicurus, \emph{Letter to
  Herodotus} 61--62, in Diogenes Laertius X). Atoms travel at equal
  speed through the void, and in aggregates they ``move in different
  directions in times so short as to be appreciable only by the reason,
  but frequently collide until the continuity of their motion is
  appreciated by sense''; the assumption that continuity persists below
  observation ``is not true in the case before us''.
\item
  \textbf{Least partless bodies} (Diodorus Cronus, reported at Sextus
  Empiricus, \emph{Pyrrhoneion Hypotyposeis} III.32). Listed in the
  doxography of material principles between the atoms of Democritus and
  Epicurus and the unjointed masses of Heraclides. Diodorus argued from
  those minima that ``a thing never is moving, but it has moved'', the
  Greek twin of Vasubandhu's \emph{na gatiḥ} three centuries earlier;
  the formula is at \emph{Adversus Mathematicos} X and is cited here
  without being read.
\item
  \textbf{The voice as spherical waves} (Vitruvius, \emph{De
  architectura} V.3.6--7,

  \begin{enumerate}
  \def\labelenumi{\alph{enumi}.}
  \setcounter{enumi}{2}
  \tightlist
  \item
    25 BCE). The voice is propelled ``by an infinite number of circles
    similar to those generated in standing water when a stone is cast
    therein'', but ``whereas the circles in water only spread
    horizontally, the voice, on the contrary, extends vertically as well
    as horizontally'', so the water-ring image is corrected into a
    spherical wave by argument from the disanalogy.
  \end{enumerate}
\item
  \textbf{The sling, applied to an orbit} (Plutarch, \emph{De facie in
  orbe lunae} 923C--D, c.~100 CE). ``The moon is saved from falling by
  its very motion and the rapidity of its revolution, just as missiles
  placed in slings are kept from falling by being whirled around in a
  circle'', with 923D adding that ``each thing is governed by its
  natural motion unless it be diverted by something else''.
\end{itemize}

Lemmas X and XI perform exactly the operation these positions dispute.
The scholia cite the ancients for what Book III concludes and leave them
silent on how Book I proceeds.

The last four entries differ in kind from the rest, and they matter most
here. The Greek and Chinese items are paradoxes about division, which
set a problem. The Indian entries answer it, and they answer it three
different ways: the arrow's flight is many motions carried by an
impression (Vaiśeṣika), or it is no motion at all because each thing
perishes where it arose (Vasubandhu), or it is possible only because a
medium stands ready to permit it (Umāsvāti). Each also carries a
stopping rule for division: the Nyāya \emph{paramāṇu} than which nothing
is smaller, the directional-parts reductio by which a touched atom would
have parts, and the Jain atom that has no space-points at all. The
Vaiśeṣika sūtras in particular state a \textbf{positive theory of
propagation}: motion is a succession of numerically distinct events,
each carried to the next by a quantity the previous one deposits, and
transmission through a medium is a chain in which what arrives is a
later member than what was sent. That is the structure the theorems
quantify. The aperture bound of §10 is protocol-universal because the
signal is a phase-space path whose total variation is unchanged by
subdivision, which is the modern form of the claim that impulsion
supplies only the first motion while the impression carries the rest. An
ancient author had therefore already framed propagation as a bounded
succession rather than a continuous traversal, and Newton's classical
preface reached for none of it.

\hypertarget{the-hierarchy-the-ancient-premises-generate}{%
\subsubsection{The hierarchy the ancient premises
generate}\label{the-hierarchy-the-ancient-premises-generate}}

Read as premises about records rather than about geometry, the cone and
the arrow are the first members of a sequence whose next member is
Newton's comparison. The apparatus of §3 covers all of them; what
changes is the order of the signal against the nuisance it must beat.

\textbf{Theorem 9.} In the model of §3, with every mark obeying
\(\delta_j\Delta_j\ge\kappa\):

\begin{enumerate}
\def\labelenumi{(\roman{enumi})}
\tightlist
\item
  \emph{The arrow.} Distinguishing rest from uniform motion at speed
  \(v\), with the initial position unknown, requires
\end{enumerate}

\[m\,v^2\tau\ \ge\ 8\,z_{1-\epsilon}^2\,\kappa,
\qquad\text{that is}\qquad \tau\,E_{\rm kin}\ \ge\ 4\,z_{1-\epsilon}^2\,\kappa .\]

\begin{enumerate}
\def\labelenumi{(\roman{enumi})}
\setcounter{enumi}{1}
\item
  \emph{Newton.} Distinguishing uniform motion from constant force, with
  the initial position and velocity unknown, requires
  \(F^2\tau^3/m\ge96z_{1-\epsilon}^2\kappa\), which is Theorem 2.
\item
  \emph{The cone.} Measuring a static taper requires nothing. With no
  recoil to propagate, the deflection grows without bound as marks
  accumulate, so a static solid has no floor.
\end{enumerate}

\emph{Proof.} For (i) the nuisance is the constant, so \(u\) ranges over
vectors with \(\sum_iu_i=0\) alone. With \(P(t)=vt\) the identity of §3
reads \(u^{\mathsf T}P=v\sum_iu_it_i=v\,T(0)\), and \(T(\tau)=0\) with
Cauchy--Schwarz gives \(|T(0)|=|\int_0^\tau T'|\le\sqrt\tau\sqrt E\).
The noise bound \(u^{\mathsf T}\Sigma u\ge2\kappa E/m\) needs only
\(\sum_iu_i=0\) and so is unchanged. Hence
\(d^2\le v^2\tau E\,m/(2\kappa E)=m v^2\tau/(2\kappa)\), and
\(d\ge2z_{1-\epsilon}\) gives the claim. For (ii) the nuisance is the
linear space, \(u^{\mathsf T}P=\frac Fm\int_0^\tau T\), and the proof of
Theorem 2 applies. For (iii) the marks leave the object unchanged, so
\(\Sigma=\operatorname{diag}(\delta_j^2)\); with \(N\) marks of
resolution \(\delta\) at each of two heights separated by \(\Delta h\),
the taper \(\vartheta\) gives
\(d^2=N\vartheta^2\Delta h^2/(2\delta^2)\), unbounded in \(N\).
\(\square\)

The two dynamical cases share one invariant. If the signal is the
\(n\)-th order departure \(P\) with \(P(0)=\dots=P^{(n-1)}(0)=0\), the
quantity the theorem bounds below is

\[m\,\bigl(P^{(n)}\bigr)^2\,\tau^{2n-1}\ \gtrsim\ \kappa ,\]

which has the dimensions of action for every \(n\): at \(n=1\) it is
\(mv^2\tau\), twice the kinetic energy times the duration, and at
\(n=2\) it is \(F^2\tau^3/m\), twice \(\tau\Delta E\). Newton's
comparison is the second rung of a ladder whose first rung is Zeno's.

\hypertarget{why-the-arrow-and-not-the-sling}{%
\subsubsection{Why the arrow and not the
sling}\label{why-the-arrow-and-not-the-sling}}

One feature of the list needs saying. Almost every ancient entry is
rectilinear or static, while the case Newton uses to \textbf{define}
centripetal force is the sling: Definition V has the stone whirled about
in a sling endeavouring to recede from the hand, and names the force
that ``retains it in its orbit'' after the planets ``perpetually drawn
aside from the rectilinear motions, which otherwise they would pursue''.

The exception is Plutarch, and it is a pointed one. \emph{De facie} 923C
has the moon kept from falling ``just as missiles placed in slings are
kept from falling by being whirled around in a circle'', which is
Newton's apparatus applied to Newton's case a millennium and a half
earlier. The inference runs the other way: for Plutarch the whirling
\emph{prevents} the fall, so speed sustains; for Newton the cord
\emph{pulls the stone inward}, and what it diverts is a straight line
that would need no cause. The same object supports opposite accounts,
and what separates them is which motion is taken to be free. Plutarch's
next clause, that ``each thing is governed by its natural motion unless
it be diverted by something else'', is as near as the passage comes,
with natural motion still doing the work inertia later does. Newton had
read him: the Classical Scholia name Plutarch among the ancients said to
have known the doctrine of gravitation, so the sling was available in a
text he was mining for authority while writing the propositions it
models.

Theorem 9 explains why the exception stayed isolated, and Aristotle
supplies the other half of the reason. \emph{De caelo} I.2 makes uniform
circular motion the natural motion of the aether, needing no cause,
which is the exact inverse of Definition V. For the heavens, then,
curved motion was the default and straight motion the anomaly, so the
sling illustrated a conclusion already held. The claim that the ancients
treat rest as natural holds below the moon and is incomplete above it;
the reference is at metadata level here. The ancient question is the
first rung, whether the thing moves at all, with position as the only
nuisance. The sling is the second rung, whether a force acts, with
uniform motion as a further nuisance, and that rung cannot be stated
until straight-line motion is held to need no account. Ancient physics
largely does not hold that, so the sling has no work to do in it. The
Vaiśeṣika comes closest by charging only the first motion to impulsion,
and stops short of making a straight continuation free and a curved one
costly.

Indian sources do supply one whirled object, and use it for a third
purpose again. The firebrand circle, \emph{alātacakra}, appears in
Vasubandhu (Pradhan 189.23--24) to argue that because contact with the
parts is successive, the awareness of a whole is really of the parts,
and at Pradhan 33.9 with \emph{āśuvṛttyā}, by rapid action. That is an
argument about sampling, and it is the ancient form of what Theorem 9
says about records: below the mesh a succession is indistinguishable
from the continuous thing it mimics. The
\href{arrow-not-sling.md}{companion note} sets this out.

\hypertarget{what-each-ancient-premise-buys}{%
\subsubsection{What each ancient premise
buys}\label{what-each-ancient-premise-buys}}

\begin{itemize}
\tightlist
\item
  \textbf{Zeno.} His conclusion holds of records: below
  \(\tau_{\rm arrow}=8z^2\kappa/(mv^2)\) every mark is consistent with
  rest. Aristotle's diagnosis explains why, since a now carries no
  record and every datum is a window.
\item
  \textbf{The Vaiśeṣika arrow and the chain of sounds.} These supply the
  positive half the Greek material lacks. Theorem 9's ladder is their
  statement made quantitative: the succession is real, each step carries
  a bounded quantity to the next, and reading any one step costs
  \(\kappa\). The sound chain is the case where an ancient author
  asserts a finite propagation time and argues it from an experiment,
  which is the premise Newton's determinate fits deny for light.
\item
  \textbf{Vasubandhu.} His conclusion is what the theorem says about
  \emph{records} rather than about the world. Below the mesh every mark
  is consistent with rest, so no recorded instant exhibits motion and
  travel is never displayed, only inferred from enough instants
  together. The theorem declines his further step from the records to
  the world, and Theorem 9(i) measures exactly how many instants are
  needed.
\item
  \textbf{Epicurus.} His is the closest ancient statement to what the
  theorems say about records. Motion presents itself as continuous to
  sense and is a succession below it, and he names the fallacy of
  extrapolating the observed continuity downward. He draws the boundary
  where sense fails; Theorem 9 draws it at a mesh computed from the mark
  cost, which is the difference between a threshold that is reported and
  one that is derived.
\item
  \textbf{Umāsvāti.} The Jain medium is the one ancient category with no
  counterpart in the theorem, and its absence is informative. Theorems 2
  and 9 need no medium: what they price is the mark, not the motion. A
  doctrine on which motion requires a permitting substance predicts
  nothing about the cost of observing it, which is the respect in which
  this entry sets a limit on how far the reconstruction reaches.
\item
  \textbf{Democritus and Chrysippus.} The dilemma about adjacent
  sections dissolves with resolution alone, and case (iii) shows it
  needs no action floor. Adjacent recorded sections are equal, and the
  inequality of the bodies is recovered cumulatively, which is
  Chrysippus's answer given a definite sense. This corrects a heuristic
  recorded earlier in this programme, that the cone is the arrow rotated
  into Euclidean time: the static problem lacks back-action, and
  back-action is the entire source of the floor.
\item
  \textbf{Newton's own objection.} Answered on his terms. The mesh
  \(\tau_*\) is a least \emph{recorded} duration and implies no least
  magnitude: it moves with \(F\), \(m\) and the confidence demanded, so
  Euclid's incommensurables survive and the quantities of the geometry
  still diminish without limit, as the Section I scholium requires. What
  is bounded is the record, at the rate \(\hbar\).
\item
  \textbf{The Mohists and Liu Hui.} They are right about records.
  Cutting stops, and the endpoint is the mesh
  \(\tau_*=(96z^2m\kappa/F^2)^{1/3}\), which is Corollary 3.
\item
  \textbf{Hui Shi.} He is right about geometry, which is Proposition 1:
  the halving never ends, and nothing in the figure resists it.
\item
  \textbf{Al-Nazzam.} The recorded trajectory below the mesh is his
  leap. The interval stays divisible, and the crossing is exhibited only
  in finitely many steps.
\item
  \textbf{The kalam time-atom.} Denied. The mesh depends on \(F\), on
  \(m\) and on the confidence demanded, so it is a dynamical resolution
  rather than a universal atom of time. A floor on action coexists with
  a continuum of instants, and Theorem 9 shows the mesh moving as the
  force changes.
\end{itemize}

The section's claim about Newton is therefore double. He looked to the
ancients for authority and found the wrong ancients, taking the
doxography of gravitation while passing over the dispute about division
that his own Lemmas turn on. And the premise that would have converted
that dispute into a theorem, an indeterminacy in the least parts of
light, is one the atomists supply in the form of a least part and one he
denied in the single place where he had a measurement.

\hypertarget{unmarked-preparations-and-prior-art}{%
\subsection{10. Unmarked preparations and prior
art}\label{unmarked-preparations-and-prior-art}}

\textbf{Preparations of unbounded extent.} Theorems 2, 5 and 6 concern
protocols that mark the trajectory. A protocol that prepares once, waits
and measures once faces a different bound. In that case the two
hypotheses differ by the phase-space displacement \((-F t^2/2m,\ Ft)\),
whose components have product exactly \(\tau\Delta E\), and
Mandelstam--Tamm in the geometric form of
\href{https://doi.org/10.1103/PhysRevLett.65.1697}{Anandan and Aharonov}
(PRL \textbf{65}, 1697, 1990), with Helstrom's discrimination bound
(\emph{Quantum Detection and Estimation Theory}, 1976), gives

\[\frac{F\tau L}{\hbar}+\frac{F\tau^2P}{2m\hbar}\ \ge\ \theta,
\qquad\tau\Delta E\ \ge\ \frac{2m\theta^2\hbar^2\tau}{(2mL+\tau P)^2},
\qquad\theta=\arcsin(1-2\epsilon),\]

for an apparatus whose position and momentum spreads are bounded by
\(L\) and \(P\). It is the integrated form of a kick bound,
\(\int_0^\tau|f(t)|L(t)\,dt\ge\hbar\theta\) with \(L(t)\) the body's
position spread at time \(t\): a force is a succession of impulses, and
an impulse is registered only through the position spread it acts on
(\href{record-distance-path-length.md}{path-length note}, Theorem K).
The floor equals \(\theta^2\hbar^2/(4LP)\) at the balanced aperture
\(L=\tau P/2m\) and falls to zero as either side grows at fixed area, so
a squeezed state of area \(\hbar/2\) evades it as a large one does. In
Newton's quantities the first inequality is
\(F\tau\cdot L+\frac{F\tau^2}{2m}\cdot P\ge\hbar\theta\), the impulse of
Proposition I against the position aperture plus the sagitta of Lemma X
against the momentum aperture. It holds for every finite adaptive
protocol of instruments, of any Kraus rank and with any apparatus
memory, because the displacement is a phase-space path whose total
variation does not grow when it is subdivided. The bound is tight: two
narrow packets separated by \(\pi\hbar/(F\tau)\) decide the comparison
with error probability of order \(\tau\Delta E/\hbar\), for arbitrarily
small \(\tau\Delta E\), at the price of an apparatus whose extent
diverges. So the floor of order \(\hbar\) requires either that the
trajectory be marked or that the laboratory be bounded, and a
non-Gaussian preparation of unbounded extent evades it.

\textbf{One accounting.} The kick bound pays for the force at the body
and Theorem 6 pays for it at the marks, and any split of the force
between the two gives a valid bound. For von Neumann marks and a force
of one sign the best split is pointwise,
\(\hbar\theta\le\int_0^\tau|f(u)|\min\bigl(L(u),R(u)\bigr)du\), where
\(R(u)=\frac1m\sum_j\Delta_jG_\tau(t_j,u)\) is the recoil length, the
marks' recoil spreads weighted by the displacement each produces at time
\(u\) relative to the chord and summed, with \(G_\tau\) the Dirichlet
Green's function of the interval: at every moment the force is paid in
the smaller of the two lengths (path-length note, Theorem M).

\textbf{Prior art.} The combination \(F^2\tau^3\gtrsim m\hbar\) is the
standard quantum limit for detecting a force on a free mass (Braginsky
and Khalili, \emph{Quantum Measurement}, 1992; the canonical statement
is Caves, Thorne, Drever, Sandberg and Zimmermann, \emph{Rev.~Mod.
Phys.} \textbf{52}, 341, 1980), and its status has been disputed since
Yuen's objection: Caves defended it (PRL \textbf{54}, 2465, 1985) and
Ozawa exhibited a measurement breaking it for free-mass position (PRL
\textbf{60}, 385, 1988). Theorem 4's coupling is von Neumann's (1932),
its inequality Robertson's (1929) and its ancestor the Heisenberg
microscope (1927); Bohr and Rosenfeld (1933) and Araki and Yanase (1960)
set limits of this kind from field measurability and from conservation
laws. That non-Gaussian preparations beat the limit, as the two-packet
preparation above does, is stated as a principle by Giovannetti, Lloyd
and Maccone (\emph{Science} \textbf{306}, 1330, 2004), and the
correlated and grid probes of §§4--5 are Yuen's contractive states and
the grid states of Gottesman, Kitaev and Preskill (PRA \textbf{64},
012310, 2001). The ingredients of §§5--6 and §10 are standard: Ozawa's
disturbance operators, Mandelstam--Tamm in the Bures angle (Uhlmann
1992; Anandan and Aharonov 1990), hybrid arguments (Bennett, Bernstein,
Brassard and Vazirani 1997), Bures-angle speed limits for processes
(Taddei et al.~2013), the position-spread form of force sensitivity
(Tsang, Wiseman and Caves 2011), and symmetry-based
information--disturbance bounds of Wigner--Araki--Yanase type. The
contribution here is universality over protocols and instruments with
explicit constants, the sharp constant of the Gaussian form, the
disturbance form of Theorem 6, the split accounting of §10, the phase
identities of §6 read in Newton's areas, and the identification of the
premise. An adversarial review of the notes behind §§5, 6 and 10 on
2026-09-23 found no false result and led to the repairs recorded in
them. The inequality \(F^2\tau^3\gtrsim m\hbar\) itself is standard.

\textbf{The classical theorem belongs to a literature it does not cite.}
Worst-case recovery of a linear functional from noisy linear data is
optimal recovery in the sense of Micchelli and Rivlin (1977), and the
apparatus series this programme built, together with the worst-case form
of §3, are results of that kind; information-based complexity (Traub,
Wasilkowski and Woźniakowski, 1988) is the same setting. The
invariant-test step of Theorem 2 is standard (Lehmann and Romano). A
referee from either community will ask why these go unmentioned, and the
honest answer is that the theorems were derived without them.

\hypertarget{obligations-before-submission}{%
\subsection{11. Obligations before
submission}\label{obligations-before-submission}}

Two, both on the history side, and both stated so that a reader can see
what the present draft rests on.

\begin{enumerate}
\def\labelenumi{\arabic{enumi}.}
\tightlist
\item
  \textbf{Historiography, to be read rather than cited.} On the
  \emph{Opticks} side, Shapiro's \emph{Fits, Passions, and Paroxysms}
  (1993) on the theory of fits and how determinate Newton meant it,
  which is where §8's central claim must be tested, and Sabra's
  \emph{Theories of Light from Descartes to Newton} (1981) for the
  context of Query 29. On the \emph{Principia} side, Guicciardini's
  \emph{Reading the Principia} (1999) and his \emph{Isaac Newton on
  Mathematical Certainty and Method} (2009) for what the limit arguments
  of §2 were for, De Gandt's \emph{Force and Geometry} (1995) directly
  on the sagitta, and Bertoloni Meli's \emph{Thinking with Objects}
  (2006) on the sling and the pendulum as objects to think with. On the
  scholia, McGuire and Rattansi, ``Newton and the `Pipes of Pan'\,''
  (1966), beside Casini and Mosley. On the ancient comparison, Sorabji's
  \emph{Time, Creation and the Continuum} (1983), which treats Zeno,
  Diodorus Cronus, Epicurean minima and kalām atoms together; von
  Rospatt (1995) on momentariness; Dhanani (1994) on kalām atomism and
  the Indian influence question; and Lloyd and Sivin's \emph{The Way and
  the Word}

  \begin{enumerate}
  \def\labelenumii{(\arabic{enumii})}
  \setcounter{enumii}{2001}
  \tightlist
  \item
    as the methodological standard that protects §9 from the charge of
    naive parallelism. Clagett (1959) for impetus, where Philoponus'
    \emph{rhopē} is the Western counterpart of \emph{saṃskāra}.
  \end{enumerate}
\item
  \textbf{Editions.} The \emph{Opticks} passages come from a
  transcription of the 1730 fourth edition and should be cited from that
  printing; the \emph{Principia} passages should be cited from Cohen and
  Whitman in place of the Motte text used here.
\end{enumerate}

Three mathematical items remain open and are not obligations of the
paper: the worst-case interval \(9\le F^2\tau_*^3/(m\kappa)\le36\) of
§3, which concerns a framework without quantum instances; whether the
minimum of the one accounting of §10 over splits and lines is attained
by some protocol, which would make it the exact floor; and whether the
constants \(1/8\) and \(1/2\) of (5) are attained, the line of best
uniform approximation being optimal for the momentum term alone.

\hypertarget{consequence-for-state}{%
\subsection{12. Consequence for STATE}\label{consequence-for-state}}

This is the synthesis the goal asked for: the Planck gap established
with Newton-age arguments and their modern equivalents, in one document
whose historical and technical halves are load-bearing for each other.
The foundations content is §§3--6: the sharp Gaussian floor, the mark
cost, the disturbance floor for every instrument and the phase that no
preparation moves. The history content is §§7--9, resting on the source
companions. The junction is §8: M3 is Robertson's inequality for
Newton's corpuscle, its universal form is that a record costs
disturbance, and Newton's fits carry a phase whose rate is proportional
to momentum and whose unit \(\Lambda p\) is a refraction invariant.
STATE's queue reduces to §11.
\end{document}
